\section{The numerical application and its exact budget}\label{sec:numerics}

\subsection{Contraction of a fixed input}
For a fixed \(s_0\in(0,1)\), define a new schedule by multiplying
every duration by \(s_0\), while preserving masks, cells and order.
Then
\[
 F'_{r,t}=s_0F_{r,t/s_0},\qquad \tau'_r=s_0\tau_r,
 \qquad A'=s_0A.
\]
Chord capacities scale by \(s_0\), and squared-chord risk scales by
\(s_0^2\). These identities follow directly by substitution in the
time integrals. They concern a new deterministic input, rather than
an unproved time-change identity for the old random process.

The adopted values from the preparation and followup records are
listed below. Fractions, rather than their decimal displays, are the
values to be checked. Their arithmetic consistency is rechecked by
the script distributed with this draft; that script does not certify
the geometric field from which the stored bounds were obtained.

\begin{center}
\begin{tabular}{@{}ll@{}}
\toprule
Quantity & Recorded value or choice\\
\midrule
Target \(c\) & \(1501/1000\)\\
Time contraction \(s_0\) & \(249/250\)\\
Scaled input mass \(A_*\) & \(5036868114751900569/3355443200000000000\)\\
Total time \(T\) & \(155699823660621/6250000000\)\\
Designated clock cap \(K\) & \(255/128\)\\
Ordinary capacity cap \(\lambda\) & \(4/5\)\\
Ordinary credit floor \(\rho\) & \(1/8\)\\
Weight exponent \(d\) & \(993/1000\)\\
Virtual exponent \(\gamma\) & \(1993/2000\)\\
Geometric and correction exponents \(\xi,\zeta\) & \(1/1000,\ 1/1000\)\\
Finite-sieve coefficient \(D_*\) & \(76/125\)\\
Spatial grid \(G\); designated directions \(g\) & \(32768;\ 16\)\\
\bottomrule
\end{tabular}
\end{center}

The recorded scaled capacity bounds and their margins are
\begin{align*}
 U_H&=\frac{159305594247}{80000000000},
 &K-U_H&=\frac{69405753}{80000000000},\\
 U_L&=\frac{2535641192953}{3200000000000},
 &\lambda-U_L&=\frac{24358807047}{3200000000000},\\
 U_{\rm tail}&=\frac{294190968417}{400000000000},
 &1-\rho-\lambda&=\frac3{40}.
\end{align*}
All margins are positive. The tail is included in the ordinary
capacity contract, not used as a substitute for the finitely many
low-height directions.

\subsection{Changing the directional coefficient without changing its meaning}
The old coefficient, used here only as a rational number, was
\[
 D_{\rm old}=\left(\sum_{k=1}^{256}\frac1{k^2}
                    +\frac1{257}+\frac1{2\cdot257^2}\right)^{-1}.
\]
For the same geometrically justified unscaled envelope integral,
the old scaled bound was
\[
 r_{\rm old}=\frac{3099987999}{3125000000},\qquad
 \kappa=d-r_{\rm old}=\frac{3137001}{3125000000}.
\]
Changing to the finite-sieve coefficient gives the inherited comparison
\begin{align}
 r_{\rm new}&=\frac{D_*}{D_{\rm old}}r_{\rm old}
       <d-\frac{7\kappa}{8},\label{eq:risktransfer}\\
 b_{\rm acc}&=d-\frac\kappa2
       =\frac{6203112999}{6250000000}=0.99249807984.
       \label{eq:bacc}
\end{align}
Thus \(b_{\rm acc}-r_{\rm new}>3\kappa/8>0\) remains for finite
subdivision and rounding. The coefficient multiplying the twelve
unscaled direction/type integrals is exactly
\[
 \frac{D_*s_0^2}{24}=\frac{392673}{15625000}.
\]
This comparison is conditional on the old bound describing the same
valid envelope. It is not a proof that every saved branch partition
is accepted by the new checker. Indeed, increasing \(D\) may make
a previously tight endpoint bound fail; the affine improvement or a
refinement must then be checked for that actual leaf.

There is a further distinction between continuum acceptance and the
discrete pair bound used by the process. If all used roots satisfy
\(\mathcal R_r\le b_{\rm acc}\), choose
\begin{equation}\label{eq:adiscrete}
 a=\frac{b_{\rm acc}+d}{2}=0.99274903992<d.
\end{equation}
Proposition~\ref{prop:pair} gives
\(\mathcal B_n^{\Gamma^c}\le b_{\rm acc}\log n+C\).
For \(\log n\ge C/(a-b_{\rm acc})\), this implies the required
\(\mathcal B_n^{\Gamma^c}\le a\log n\).
One must not silently absorb the additive error while retaining the
same coefficient \(b_{\rm acc}\).

\subsection{Moment choices and the fixed loss budget}
The exponent conditions hold because
\[
 1-d-2\xi-\zeta=\frac1{250}>0,\qquad
 \gamma-d=\frac7{2000}>0.
\]
For example \(m=4096\) makes the exponents in
\eqref{eq:etatail} and \eqref{eq:virtualtail}, before point-law
transport, equal to \(-13.377\) and \(-12.328\).
These are illustrative fixed choices, not a uniform statement about
the hidden constants for every order.

For the tag bound, \(C_{\rm line}=4G=131072\). The recorded fixed
budget \(\varepsilon=10^{-5}\), with one third assigned to tags,
can use
\[
 e=\frac{1-\rho-\lambda}{4TC_{\rm line}},\qquad
 b=348294609696,\quad k=116,\quad
 m_{\rm tag}=bk+3=40402174724739.
\]
The rational witnesses
\(b(e/2)\ge1\) and
\(2^{-k}\le (1/300000)e^2/(64C_{\rm line})\) imply the desired
small main term in \eqref{eq:tagprob}, by Bernoulli's inequality
\((1+e/2)^b\ge2\). Its enormous but fixed degree is permitted by
the proof: a general fixed-order moment argument is used before
substituting this integer. No such power is expanded into individual
terms in the arithmetic check. This choice supplies neither a useful
explicit \(n_0\) nor an estimate for practical simulation size.

\begin{corollary}[Candidate application, conditional on full input checking]\label{cor:target}
Suppose the candidate accepted schedule has mass \(A_*\), the stated
capacity bounds, and a valid geometric certificate giving risk at
most \(b_{\rm acc}\) at every used root's deadline. Suppose the
accepted unions, durations and reference prefixes have the matching
specified in Proposition~\ref{prop:contract}. Then the analytic
reduction yields \eqref{eq:target}.
\end{corollary}
\begin{proof}
The finite sieve, positive margins and \eqref{eq:adiscrete} supply
the geometric hypotheses; the displayed parameters satisfy
\eqref{eq:exponents}. Theorem~\ref{thm:bridge} applies.
Alternatively, retain just one error budget: choose \(n\) large
enough that \(A_n\ge A_*-\Delta/4\) and the total early-retirement
point-law probability is at most \(10^{-5}\). Proposition~\ref{prop:loss}
and the strict inequality in Section~\ref{sec:overview} give
\(\E|S_n|>cn\). A legal realization then satisfies
\(1501n\le1000|S_n|\).
\end{proof}

A separately verified pruning loss \(P\le\Delta/4\) can also be
paid in the same ledger, provided the resulting schedule continues
to satisfy the required geometric and matching conditions. With
grid error \(\Delta/4\),
\[
 (1-10^{-5})(A_* -\Delta/4-\Delta/4)-c
 =0.000036826382327236234293878078460693359375>0.
\]
The two occurrences of \(\Delta/4\) pay different costs; neither
is reused. This is a budget lemma, not a claim that a pruned input
has actually been constructed or checked for this manuscript.

\section{Verification status, correspondence and provenance}\label{sec:status}

\subsection{What this document adds}
This document is an English mathematical exposition of the preparation
notes and their two followups~\cite{preparation,followup,second},
organized around the exact loss identity.
It adopts finite inclusion-exclusion in place of the older infinite
Mobius argument, affine rational integral bounds in place of mandatory
logarithmic integration, and the linewise cell count in the tag proof.
The arguments are written out for a reader, rather than replaced by
Lean declarations. They remain a research proof draft under review.
Editing and arithmetic checks are not independent verification of the
entire stochastic argument or the concrete numerical certificate.

The paper is self-contained in the sense that its mathematical
definitions, proof steps and the checking contract are stated here.
It is not a standalone accepted numerical certificate: the large
reference field and all geometric leaves are not embedded in the
text or in the small companion data file. The numerical corollary
retains those explicit premises.

\subsection{The fixed Lean snapshot}
The comparison here is with the public snapshot~\cite{formal} whose clean-history
commit is \texttt{282b636}, based on the audit of 10 November 2026,
00:22:46 JST. It contains 3,242 original Lean modules and reports
839,734 source lines. This count measures source inventory, not the
percentage of the final theorem proved. The fixed snapshot records:
\begin{itemize}
\item actual one-phase exponential estimates and uniform power tails;
\item actual frozen-schedule clock-cause smallness;
\item an output theorem combining the five causes under explicit
concrete-input hypotheses;
\item 99 of 112 low-height capacity rows, and the specified nonaxis
capacity table with 28 of 28 rows;
\item 20 of 80 canonical leaves at the single benchmark phase 15301,
with 80 of 320 symmetry images.
\end{itemize}
The last count is not 25 percent of the complete geometric input.
The full input has 15,348 phases and 155,402 leaf records.

The principal formal endpoint is
\begin{center}
\small\path{FrozenCapacity.exists_uniform_accepted_frozen_output_of_inputs}
\end{center}
in \path{NoThreeLine/FrozenFiveCauseOutput.lean}.
Its hypotheses include all 112 accepted rows; a fixed number of
checked accepted phases; their matching to the frozen reference with
strictly increasing indices; their duration bounds and fixed total
time; actual grid input mass above \(1501/1000\); and the actual
normalized pair-sum bound at every used root's deadline, with fixed
\(0\le a<d<\gamma<1\) and \(0<\zeta<1-d\).
It does not assert that these inputs have been constructed for all
sufficiently large \(n\).

Two intermediate entry points are
\path{NoThreeLine/ProtectedCellOnePhasePowerTail.lean} and
\path{NoThreeLine/FrozenClockCauseSmall.lean}.
The general expectation-to-output interface is in
\path{NoThreeLine/ScheduleFiveCauseOutput.lean}.
The paper's explanatory grouping is not a claim that every displayed
lemma is a verbatim translation of one Lean declaration. The formal
source, with its actual hypotheses, determines the formal scope.

\subsection{The remaining numerical and mathematical obligations}
To turn Corollary~\ref{cor:target} into an unconditional claim, one
must complete a single compatible chain of evidence:
\begin{enumerate}
\item Certify every adopted root region with its correct field,
parity labels, angular coverage and symmetry images. If using the
new finite-sieve coefficient, establish acceptance under that
coefficient, not just under the older one.
\item Establish chronological phase matching and all cumulative
reference-prefix bounds. A local accepted rectangle alone is not
a valid complete schedule.
\item Check union mass, capacity maxima, axis strips and direction
tails for that same input. Complete the remaining finite row
obligations if using the frozen Lean interface.
\item Transfer those continuum bounds to the actual grid input and
actual root-deadline pair sum with one sufficiently large threshold
uniform in all used roots. The additive-error step in
\eqref{eq:adiscrete} is part of this obligation.
\end{enumerate}
The probabilistic proof also deserves direct mathematical examination,
especially common stopping of mixed hierarchies, whole-event energy,
conditional restart at the first line birth, and the exceptional
configurations in the short-star expansion. The recorded AI reviews
and partial formalization are evidence at their stated scopes; they
do not convert the remaining premises into facts.

\subsection{Companion files and what was checked for this edition}
The editable source is \path{main.tex} with its section files.
The small file \path{data/parameters.json} records the exact numerical
choices and the finite-sieve factors. Running
\path{scripts/check_arithmetic.py} checks the mass and capacity
margins, the coefficient transfer, the discrete-risk slack,
the two fixed loss budgets, the finite product and coprimality
conditions, and the tag-order witnesses. Its report is a verification
of those finite arithmetic statements, not an execution of the full
geometric checker or a check of the whole proof.

\path{SOURCE_MAP.json} identifies the source notes and the frozen
public modules used in preparing this edition. Source identifiers
and hashes serve reproducibility; their role is distinct from the
soundness arguments in Section~\ref{sec:certificates}.
The PDF has been typeset and visually checked. This edition performs
no new large angular computation and no full Lean rebuild.

\subsection{AI use and licensing}
This work has involved substantial use of OpenAI ChatGPT/Codex and
AI reviews supplied by the project maintainer. AI systems contributed
mathematical proposals and derivations, numerical programs, exact
arithmetic checks, Lean definitions and proofs, build orchestration,
review, translation and exposition. Separate AI contexts reviewed
parts of the earlier work. This English edition was assembled and
checked by the assistant in this task; no new independent human
review or full independent proof audit is claimed.

The human maintainer chose the direction and requested experiments
and revisions as a hobby project. This disclosure does not attribute
every mathematical or programming step to human authorship. No named
author, affiliation or personal contact address is supplied. The
repository's hosting account is public. The original project material
and this draft are distributed under the MIT license supplied in
\path{LICENSE}; retain its notices when redistributing. No claim of
anonymous hosting or of formal completion is made.

\begin{thebibliography}{9}
\bibitem{preparation}
Project preparation notes, \emph{Integrated proof draft preceding
Lean formalization}, 8--9 November 2026 (Japanese).
\href{https://github.com/agolkjaoi/no-three-in-line-lean/blob/282b636c160d049e3ef4d82f619fd44d98f8cf1d/docs/research-notes/preparation/PROOF_FOR_LEAN_JA.txt}{Public fixed-version text}.
This draft uses its process, probability bridge, uniform geometry,
and stopping/risk arguments.
\bibitem{followup}
Project first followup, 9 November 2026 (Japanese).
\href{https://github.com/agolkjaoi/no-three-in-line-lean/blob/282b636c160d049e3ef4d82f619fd44d98f8cf1d/docs/research-notes/first-followup/FOLLOWUP_FOR_REVIEW_JA.txt}{Public fixed-version text}.
Source of the remaining-exposure identity, fixed loss budget,
finite height sum and strengthened capacity contract.
\bibitem{second}
Project second followup, 9 November 2026 (Japanese).
\href{https://github.com/agolkjaoi/no-three-in-line-lean/blob/282b636c160d049e3ef4d82f619fd44d98f8cf1d/docs/research-notes/second-followup/SECOND_FOLLOWUP_FOR_REVIEW_JA.txt}{Public fixed-version text}.
Source of the finite sieve, affine integral bounds and linewise
cell count.
\bibitem{formal}
\emph{no-three-in-line-lean}, partial formalization snapshot,
10 November 2026.
\href{https://github.com/agolkjaoi/no-three-in-line-lean/tree/282b636c160d049e3ef4d82f619fd44d98f8cf1d}{Source snapshot};
\href{https://github.com/agolkjaoi/no-three-in-line-lean/blob/282b636c160d049e3ef4d82f619fd44d98f8cf1d/STATUS.md}{status and remaining obligations}.
\end{thebibliography}

